\documentclass[12pt]{article}

\usepackage[utf8]{inputenc}
\usepackage[T1]{fontenc}
\usepackage[brazil]{babel}

\usepackage{amsmath, amssymb, amsthm}
\usepackage{mathtools}
\usepackage{geometry}
\geometry{margin=2.5cm}

\usepackage{setspace}
\onehalfspacing

\usepackage{hyperref}
\usepackage{csquotes}

\usepackage[
    backend=biber,
    style=authoryear
]{biblatex}

\addbibresource{/mnt/c/Users/nicho/Documents/nicholasvoltani.github.io/bibliography.bib}

\begin{document}

\hypertarget{conceitos-introdutuxf3rios}{%
\section{Conceitos introdutórios}\label{conceitos-introdutuxf3rios}}

\hypertarget{conjuntos-parcialmentetotalmente-ordenados}{%
\subsection{Conjuntos parcialmente/totalmente
ordenados}\label{conjuntos-parcialmentetotalmente-ordenados}}

Um conjunto (não-vazio) \(X\) é dito \emph{parcialmente ordenado} se ele
possui uma relação \(\leq\) que satisfaça\footnote{Note que \(\forall\)
  é um ``A'' invertido, que em inglês se lê ``\emph{for \textbf{a}ll}''.
  Por exemplo, a asserção \(\forall x \in X: x \leq x\) lê: ``Para todo
  \(x\) pertencente ao conjunto \(X\), vale que \(x \leq x\)''. Ademais,
  caso valha a asserção lógica \(A \land B\), lê-se: ``\(A\) \textbf{e}
  \(B\) são verdadeiros''.} - \emph{Reflexividade}:
\(\forall x \in X: x \leq x\) - \emph{Transitividade}:
\(\forall x, y, z \in X: (x \leq y) \land (y \leq z) \implies x \leq z\)
- \emph{Antissimetria}:
\(\forall x, y \in X: (x\leq y) \land (y \leq x) \implies x = y\)

Essa relação \(\geq\) é dita ser uma \emph{relação de ordem parcial} (ou
relação \emph{parcial} de ordem).

Caso \emph{todos} os elementos em \(X\) sejam ``comparáveis'' através de
\(\leq\) --- ou seja,\footnote{Caso valha a asserção lógica
  \(A \lor B\), lê-se: ``\(A\) \textbf{ou} \(B\) é/são verdadeiro(s)''.
  Note que este ``ou'' é \emph{inclusivo}, ou seja, quer dizer ``\(A\) é
  verdadeiro, \(B\) é verdadeiro, ou \(A\) \textbf{e} \(B\) são
  verdadeiros''.} \[
\forall x, y \in X: (x \leq y) \lor (y \leq x)
\] dizemos que o conjunto \(X\) é \emph{totalmente ordenado}, e que
\(\leq\) é uma relação de ordem \emph{total} (ou relação \emph{total} de
ordem).

Note que, naturalmente, nem todas as relações parciais de ordem são
totais. Seja \(X = \{ 1,2,3\}\), e
\(\mathcal{P}(X) = \{ \emptyset, \{ 1 \}, \{ 2 \}, \{ 3 \}, \{ 1,2 \}, \{ 1,3 \}, \{ 2,3 \}, X \}\)
o conjunto das partes de \(X\) (em inglês: \emph{power set}). Verifique
que \(\mathcal{P}(X)\), com a relação de inclusão de conjuntos
\(\subseteq\), é uma relação \emph{parcial} de ordem, mas que \emph{não
é total} --- ou seja, existem elementos em \(\mathcal{P}(X)\) que não
são comparáveis sob \(\subseteq\).

(Ao leitor pedante: trivialmente pode-se definir \(x \preceq y\) --- a
operação ``para o outro lado'' --- como \(y \succeq x\).)

\hypertarget{preferuxeancias-racionais}{%
\subsection{Preferências racionais}\label{preferuxeancias-racionais}}

(\textbf{Mas-Colell1995?}) definem \emph{preferências racionais}
\(\succeq\) como ``relações \emph{completas}
(\(\forall x,y: x \succeq y \lor y \succeq x\)) e \emph{transitivas}''
--- ou seja, vide acima, são relações de ordem \emph{total} sobre um
espaço de cestas de bens.

Vide definição \textbf{1.B.2.} em (\textbf{Mas-Colell1995?}), uma função
\(u: X \to \mathbb{R}\) é uma \emph{função utilidade que representa uma
relação de ordem} \(\succeq\) caso \[
\forall x, y \in X: x \succeq y \iff u(x) \geq u(y)
\] (naturalmente com a relação de ordem \(\geq\) usual dos números reais
\(\mathbb{R}\).)

A relação de ordem sobre os reais \(\mathbb{R}\) \emph{é uma relação
total}, pois sabemos que, para quaisquer números reais
\(x, y \in \mathbb{R}\), vale que ou \(x>y\), ou \(x<y\), ou \(x=y\);
dito de outra forma, sempre vale que ou \(x \geq y\), ou \(y\geq x\).

Dessa forma, caso uma relação de preferências \(\succeq\), sobre um
conjunto \(X\) de cestas de bens, seja representável por uma função
utilidade \(u:X \to \mathbb{R}\), então \emph{necessariamente} esta
relação tem de ser \emph{racional} (i.e.~relação \emph{total} de ordem).
A demonstração da proposição \textbf{1.B.2} em
(\textbf{Mas-Colell1995?}) segue o seguinte esquema: para quaisquer
\(x, y \in X\), como \(u(x)\) e \(u(y)\) estão sobre um conjunto
totalmente ordenado \(\mathbb{R}\), então vale que ou
\(u(x) \geq u(y)\), ou \(u(y) \geq u(x)\), e como pressupõe-se que
\(\succeq\) é representada por \(u\), então conclui-se que ou
\(x \succeq y\) ou \(y \succeq x\) --- portanto, esta relação
\(\succeq\) é total.

Ou seja, caso uma relação de preferências seja representável por uma
função utilidade, então ela será racional. A recíproca não é verdadeira:
há preferências racionais que não são representáveis por funções
utilidade, uma das quais é a ordem lexicográfica (exercício 1.5 abaixo).

\%\% --- \%\% \# Exercícios \#\# Relações de ordem e preferências
racionais

(\textbf{Exercício 1.1}) Demonstre que, a partir de uma relação de ordem
\(\succeq\) em um espaço \(X\), é possível definir uma relação \(\sim\)
que \[x
\forall x, y \in X: x \sim y \iff (x \succeq y)\,  \land  \, (y \succeq x)
\] (Tradução: ``Para todo \(x\) e \(y\) pertencentes a \(X\), vale que
\(x \sim y\) se, e somente se, valer que \(x \succeq y\) e também
\(y \succeq x\).)

Essa relação \(\sim\) satisfaz as propriedades: 1. \emph{Reflexividade}:
\(\forall x \in X: x \sim x\) 2. \emph{Comutatividade}:
\(\forall x, y \in X: x \sim y \iff y \sim x\) 3. \emph{Transitividade}:
\(\forall x, y, z \in X: (x \sim y) \land y( \sim z) \implies x \sim z\)

Por satisfazer estas propriedades, \(\sim\) é dita uma \emph{relação de
equivalência}.\footnote{Note que a relação de igualdade \(x=y\) satisfaz
  essas propriedades. Relações de equivalência são mais gerais, e menos
  restritivas, do que relações de igualdade.}

(\textbf{Exercício 1.2}) Dada uma relação de equivalência \(\sim\) sobre
algum conjunto (não-vazio) \(X\), então é possível definir, para cada
\(x \in X\), sua \emph{classe de equivalência} \([x]\) como sendo o
conjunto\footnote{Tradução matemática: \([x]\) \emph{é definido} como
  ``o conjunto de elementos \(y \in X\) tais que \(x \sim y\)''. Note
  que estes elementos \(y\) são um mero índice, i.e.~não importa qual
  letra seja usada; poderia igualmente ser escrito como ``o conjunto de
  elementos \(z \in X\)'', ou ``\ldots elementos \(\xi \in X\)''. O
  único elemento que tem alguma ``concretude'' aqui (além do dado
  conjunto \(X\)) é um dado elemento \(x\), que foi pré-determinado no
  enunciado.} \[
[x] \coloneqq \{y \in X: y \sim x\}
\]

Ou seja, é o conjunto de elementos em \(X\) que são equivalentes (sob
\(\sim\)) a \(x\) (note que \(x \in [x]\) --- por quê?).

Demonstre que classes de equivalência são \textbf{disjuntas}: se há
algum elemento em comum entre duas classes de equivalência \([x]\) e
\([y]\), então elas são a mesma classe de equivalência. Dito de outra
forma: se \(x \nsim y\) (ou seja, que \textbf{não} seja verdade que
\(x \sim y\)), então não haverá interseção entre suas classes de
equivalência --- i.e.~\([x] \cap [y] = \emptyset\) ---, e, portanto,
\([x] \neq [y]\). Escrito em ``matematiquês'', esta proposição
lê:\footnote{\(A \implies B\) lê ``\(A\) implica \(B\)'', ou seja, \(A\)
  ser verdadeiro \emph{implica que} \(B\) também seja verdadeiro.
  \(A \iff B\) quer dizer que \(A\) é verdadeiro \emph{se e somente se}
  \(B\) for verdadeiro; é o mesmo que \(A \implies B\) \textbf{e também}
  \(B \implies A\).} \[
\forall x, y \in X: x \nsim y \iff [x] \neq [y]
\]

A importância deste exercício é que, dada uma relação de preferência
racional \(\succeq\) sobre algum conjunto de cestas de bens \(X\), é
possível definir uma relação de equivalência \(\sim\) , a qual descreve
uma ``relação de indiferença'' entre cestas de \(X\). As \emph{curvas de
indiferença} são, portanto, classes de equivalência sob esta relação de
equivalência induzida por esta preferência. Segue do resultado acima,
portanto, que \emph{curvas de indiferença diferentes não se cruzam}!

(\textbf{Exercício 1.3}) Uma função monotônica \(f:X \to Y\) (definida
sobre conjuntos parcialmente ordenados \((X, \succeq_{X})\) e
\((Y, \succeq_{Y})\)) é uma função que \emph{preserva relações de
ordem}. Ou seja, se temos \(x \succeq_{X} y\) em \(X\), então
\(f(x) \succeq_{Y} f(y)\) em \(Y\). Em matematiquês: \[
\forall x, y \in X: x \succeq_{X} y \implies f(x) \succeq_{Y} f(y)
\]

Mostre que relações de equivalência (induzidas por relações de ordem)
são também preservadas por \(f\): se \(x \sim y\), então
\(f(x) \sim f(y)\).

Isso quer dizer que, embora ``os valores mudem'' por conta da função
\(f\), \emph{o formato das curvas de indiferença permanecem os mesmos}.

\textbf{Obs}: É comum se dizer que funções monotônicas são ``funções
não-decrescentes'', e que funções \emph{estritamente} monotônicas são
``funções \emph{estritamente} crescentes'' --- isso, porém, segue da
definição de acima\footnote{Usualmente também assume-se que o domínio da
  função \(f\) seja \emph{totalmente} ordenado, que é o caso de
  \(\mathbb{R}\): todos os elementos são comparáveis entre si sob a
  relação de ordem. Mas isso é mero pedantismo: se dois elementos
  \(x, y \in X\) não fossem comparáveis sob relação de ordem, então já
  não haveria nada entre eles para ser preservado por \(f\), e tudo se
  passaria como se nada tivesse se passado!}. Funções monotônicas em
geral preservam \(\succeq\), mas pode haver casos em que, por exemplo,
\(x \succ_{X} y\) \textbf{e} \(\mathbf{x \nsim_{X} y}\), mas
\(f(x) \succeq_{Y} f(y)\), ou até \(f(x) \sim_{Y} f(y)\) --- por
exemplo, funções em \(\mathbb{R}\) que sejam crescentes mas que possuam
``platôs'' nos quais certos intervalos estejam ``na mesma altura'' num
gráfico; são, portanto, funções ``não-decrescentes''. Por outro lado,
funções \emph{estritamente} monotônicas preservam somente \(\succ\) ---
se \(x \succ_{X} y\), então garante-se que \(f(x) \succ_{Y} f(y)\) ---
e, portanto, são funções \emph{estritamente} crescentes.

(\textbf{Exercício 1.4}) Considere uma relação de preferência definida
sobre \(\mathbb{R}^2\) definida como \[
(x_{1},x_{2}) \succ (y_{1}, y_{2}) \iff x_{1}+x_{2} > y_{1}+y_{2}
\]

Essa relação de preferência satisfaz não-saciedade local?

(\textbf{Exercício 1.5}) A ordem lexicográfica \(\succeq\) é uma relação
de ordem definida para quaisquer \(x=(x_{1},x_{2}) \in \mathbb{R}^{2}\)
e \(y =(y_{1},y_{2})\in \mathbb{R}^{2}\) como \[
x \succeq y \iff (x_{1} \geq y_{1}) \lor [(x_{1} = y_{1}) \land (x_{2} \geq y_{2})]
\] É chamada ``lexicográfica'', pois é um ordenamento alfabético: por
exemplo, \(ab \succeq aa\), pois \(a=a\) (ou seja, \(x_{1}=y_{1}\)) e
\(b\geq a\) (ou seja, \(x_{2}\geq y_{2}\)).

Em outras notícias, o velho Jorge é alcoólatra. Seus amigos falam em voz
baixa sobre suas curvas de nível lexicográficas: ele sempre prefere
estritamente a cesta de consumo que contém a maior quantidade de bebida,
independentemente da quantidade dos outros bens na cesta; se duas cestas
tiverem a mesma quantidade de bebida, ele prefere estritamente a cesta
que tem a maior quantidade dos outros bens. Esboce as preferências do
velho Jorge em um diagrama. Qual dos axiomas (completude e
transitividade) é violado por tal ordenação? Ela é uma relação de ordem
\emph{contínua}?

\%\% --- \%\% \#\#\# Referências - MAS-COLELL, Andreu; WHINSTON, Michael
Dennis; GREEN, Jerry R. \textbf{Microeconomic theory}. New York: Oxford
University Press, 1995.

\printbibliography

\end{document}
